The following meta-review was prepared by the program committee for the 2025 IEEE Symposium on Security and Privacy (S\&P) as part of the review process as detailed in the call for papers.

\subsection{Summary} This paper presents a novel method to verify computations performed using Fully Homomorphic Encryption (FHE), named DataSeal. DataSeal integrates Algorithm-Based Fault Tolerance (ABFT) into FHE to ensure the integrity of computations with low overhead. The method is designed for matrix operations and neural network inference tasks, leveraging secret client-specific checksums for verification. The proposed system is benchmarked against existing solutions like MAC, ZKP, and TEE, showing superior efficiency as the problem size increases.

\subsection{Scientific Contributions} \begin{itemize} \item Creates a New Tool to Enable Future Science \item Provides a Valuable Step Forward in an Established Field \end{itemize}

\subsection{Reasons for Acceptance} \begin{enumerate} \item The verifiability of computations in FHE has been a significant challenge. DataSeal offers an effective solution that addresses this gap with minimal overhead.

\item DataSeal opens new opportunities for researchers working with FHE and matrix operations. Its lightweight design makes it a valuable tool for future applications in secure computing.

\item The work builds on existing methods (ABFT and FHE) but provides significant advancements by integrating these techniques in a new way to enable verifiable encrypted computations. \end{enumerate}

\subsection {Noteworthy Concerns} The authors need to expand their evaluation to other computational workloads and real-world attack scenarios to strengthen the generalizability and robustness of DataSeal.